\begin{table}[ht]
\caption{\textbf{Example of Gemini reasoning on \sustainable}}
\centering
% \resizebox{\textwidth}{!}{
\begin{tabular}{p{13cm}}
\toprule
% \hspace{5.5cm} \sustainable \\
% \midrule
\textit{\textbf{Query}} \\
\midrule
How good is it to reuse water from plant pots?

I'm living in an apartment, and after I water my plants the water goes to plates below the pots. The pots are in a metallic structure above the plates, so I can take the plates to reuse the water (throwing it at the plants again).

This reuse seems beneficial, because I think I can get rid of mosquitoes that would reproduce in the stagnated water. And also some nutrients of the soil (as well as earthworms) can return to the vase.

Is there some negative points in doing that? \\
\midrule
\textit{\textbf{LLM Reasoning steps}} \\
\midrule
\#\# Plant Pot Water Reuse \\
\\
**(1) Essential Problem:** \\
\\
The user wants to know if it is beneficial or detrimental to reuse the water collected from plant pot drainage plates. \\
They are concerned about potential mosquito breeding and the impact on plant growth. \\
\\
**(2) Reasoning:** \\
\\
**Reused Water Composition:**\\
\\
- Potentially contains dissolved nutrients, minerals and salts from the soil. \\
- May contain microorganisms. \\
- Could harbor mosquito larvae if stagnant. \\
\\
**Negative Impact on Plant Growth:** \\
\\
- Stagnant water can attract mosquitoes. \\
- High salt concentrations from fertilizers can harm plants. \\
- Waterborne pathogens could potentially harm plants. \\
\\
**(3) Draft Answer:** \\
\\
\#\# Reusing Water from Plant Pots: A Balancing Act \\
\\
It's commendable that you're looking for ways to conserve water and potentially benefit your plants. \\
Reusing water from plant pot drainage plates can be a good idea, but there are some drawbacks to consider: \\
\\
* **Mosquito Breeding:** Stagnant water in the plates can become a breeding ground for mosquitoes, \\
posing a health risk. Regularly empty and clean the plates to prevent this. \\
* **Salt Accumulation:** Repeatedly using the same water can lead to a buildup of salts from fertilizers, which can harm plants. \\
* **Pathogen Risk:** While less common, waterborne pathogens could potentially harm your plants. \\
Using filtered or boiled water can minimize this risk. \\
\\
% **Recommendations:** \\
% \\
% * **Empty and clean the drainage plates regularly** to prevent mosquito breeding. \\
% * **Dilute the reused water with fresh water** to avoid salt buildup. \\
% * **Consider using filtered or boiled water** to minimize the risk of pathogens. \\
% * **Monitor your plants** for any signs of distress or nutrient deficiency. \\
% \\
% By taking these precautions, you can reap the benefits of reusing water while minimizing potential risks. \\
% Remember, it's always best to observe your plants and adjust your practices accordingly. \\
\bottomrule
\end{tabular}

% \begin{tabular}{p{13cm}}
% \toprule
% \textit{\textbf{Example positive document}} \\
% \midrule
% Soluble Salts

% Soluble salts may accumulate on the top of the soil, forming a yellow or white crust. A ring of salt deposits may form around the pot at the soil line or around the drainage hole. Salts may also build up on the outside of clay pots. In house plants, signs of excess soluble salts include reduced growth, brown leaf tips, dropping of lower leaves, small new growth, dead root tips, and wilting.

% Soluble salts are minerals dissolved in water. Fertilizer dissolved in water becomes a soluble salt. When water evaporates from the soil, the minerals or salts stay behind. As the salts in the soil become more and more concentrated, it becomes more difficult for plants to take up water. If salts build up to an extremely high level, water can be taken out of the root tips, causing them to die. High levels of soluble salts damage the roots directly, weakening the plant and making it more susceptible to attack from insects and diseases. One of the most common problems associated with high salt levels is root rot.

% The best way to prevent soluble salt injury is to stop the salts from building up. When watering, allow some water to drain through the container and then empty the saucer. Do not allow the pot to sit in water. If the drained water is absorbed by the soil, the salts that were washed out are reabsorbed through the drainage hole or directly through a clay pot. 

% ...\\

% \bottomrule
% \end{tabular}
% }
\label{tab:example_sustainable_living_reason}
\end{table}